\chapter{РАЗРАБОТКА АЛГОРИТМОВ ВНУТРИСХЕМНОГО ИЗМЕРЕНИЯ ЧАСТОТНЫХ ХАРАКТЕРИСТИК}

Система управления вторичного источника электропитания представлена в виде обобщенной системы аналого-цифрового преобразования и цифровой обработки сигналов с воздействием на объект управления посредством импульсной модуляции с постоянным периодом ${T_{d}}$. (рисунок~\ref{fig:discrete_time_system}). Период ${T_{d}}$ выполнения операций в системе соответствует базовой частоте дискретного регулирования ${f_{d} = 1 / T_{d}}$.

\begin{figure}[H]
	\centering
	\includegraphics [width=0.85\linewidth] {frequencyResponse/discreteTimeSystem.drawio.png}

	\caption{Рассматриваемая система}
	\label{fig:discrete_time_system}
\end{figure}

В канал управления системы аддитивно инжектируется сигнал в виде дискретной последовательности ${z[n]}$. Откликом системы являются два континуальных сигнала: ${x_{c}(t)}$, ${y_{c}(t)}$. Система преобразует континуальные сигналы ${x_{c}(t)}$, ${y_{c}(t)}$ в дискретные последовательности ${x[n]}$, ${y[n]}$. Преобразование осуществляется посредством дискретизации:
\begin{equation}\label{eqn:sampleRule}
    \begin{split}
        &x[n] = x_{c} (n T_s), \\
        &y[n] = y_{c} (n T_s),
    \end{split}
\end{equation}
где ${T_{s}}$ -- период дискретизации, ${T_{s} \geq T_{d}}$.

Во время измерения отклики ${x[n]}$, ${y[n]}$ хранятся в буфере, выделенном под внутрисхемное измерение в памяти встроенной системы управления. Размер буфера ограничен ${2N_\text{max}}$-значениями.

Коэффициенты разложения дискретных последовательностей
\begin{equation*}
    \begin{aligned}
        &x\left[n\right] = \cfrac{1}{N} \sum_{l=0}^{N - 1} {X\left[l\right] e^{j\cfrac{2\pi}{N}ln}}, \\
        &y\left[n\right] = \cfrac{1}{N} \sum_{l=0}^{N - 1} {Y\left[l\right] e^{j\cfrac{2\pi}{N}ln}}
    \end{aligned}
\end{equation*}
определяются по $N$-последовательным отсчетам дискретных последовательностей ${x[n]}$, ${y[n]}$:
\begin{equation}\label{eqn:DFT}
    \begin{aligned}
        &X\left[l\right] = \sum_{n=0}^{N - 1} {x\left[n\right] e^{-j\cfrac{2\pi}{N}nl}}, \\
        &Y\left[l\right] = \sum_{n=0}^{N - 1} {y\left[n\right] e^{-j\cfrac{2\pi}{N}nl}}.
    \end{aligned}
\end{equation}

Коэффициенты ${X\left[l\right], Y\left[l\right] \in \mathbb{C}}$, представляют собой дискретные периодические последовательности с периодом $N$.

При условии, что измерительное окно длительностью ${T_{h}}$ вмещает целое количество $N$ периодов дискретизации ${T_{s}}$ (${T_{h} = N \mathord{\cdot} T_{s}}$) дискретные последовательности ${X\left[l\right]}$, ${Y\left[l\right]}$ на отрезке ${l \in \left[0; N - 1 \right]}$ длины $N$ содержат коэффициенты разложения ${x\left[n\right]}$, ${y\left[n\right]}$ в диапазоне частот ${f \in \left[0, f_{s} / 2\right]}$ с шагом по частоте ${\Delta f=1 / T_{h}}$.

Составляющим с частотой ${f_{h}}$ соответствуют ${X\left[l\right]}$, ${Y\left[l\right]}$ при ${l = 1}$, а с частотой ${f_{s} / 2}$ -- при ${l = \bigl \lfloor (N - 1)/2 \bigr \rfloor}$.

В работе рассматриваются алгоритмы измерения частотных характеристик, основанные на воздействии на систему дискретной последовательностью ${z[n]}$ моногармонического типа с амплитудой ${Z}$ и частотой ${\omega_z = 2\pi f_z}$:
\begin{equation*}
    z[n] = Z \cos{(n \omega_z T_{d})}.
\end{equation*}
В ходе измерения период инжектируемого сигнала ${T_z = 1 / f_z}$, а, в общем случае, и амплитуда $Z$, последовательно изменяются. В серии из $K$-измерений каждому $k$-му измерению, ${k \in \{1, \dots, K\}}$, соответствует пара параметров ${T^{(k)}_{z}}$, ${Z^{(k)}}$, следовательно, и набор дискретных последовательностей ${x^{(k)}[n]}$, ${y^{(k)}[n]}$, а также коэффициентов ${X^{(k)}\left[l\right]}$, ${Y^{(k)}\left[l\right]}$.

Под частотной характеристикой понимается набор из $K$-пар значений ${f^{(k)}_{z}}$, ${H^{(k)}}$, где в каждом $k$-ом измерении
\begin{equation}\label{eqn:complexFreqResp}
    H^{(k)} = \cfrac{Y^{(k)}\left[L^{(k)}\right]}{X^{(k)}\left[L^{(k)}\right]} = \left| H^{(k)} \right| e^{j \angle H^{(k)}},
\end{equation}
${L^{(k)} = {T^{(k)}_{h}} / {T^{(k)}_{z}}}$.

Пары значений ${f^{(k)}_{z}}$, ${\left| H^{(k)} \right|}$ и ${f^{(k)}_{z}}$, ${\angle H^{(k)}}$ образуют амплитудно-частотную и фазо-частотную характеристики, соответственно.
\begin{equation}\label{eqn:abs_angleFreqResp}
    \begin{aligned}
        &\left| H^{(k)} \right| = \cfrac{\left| Y^{(k)}\left[L^{(k)}\right] \right| }{ \left| X^{(k)}\left[L^{(k)}\right] \right|}, \\
        &\angle H^{(k)} = \angle Y^{(k)}\left[L^{(k)}\right] - \angle X^{(k)}\left[L^{(k)}\right].
    \end{aligned}
\end{equation}

\section{Алгоритм с изменяемой длительностью измерительного окна, равной периоду инжектируемого сигнала}\label{sec:N_var}

Данный алгоритм характеризуется тем, что в серии из $K$ изменений периода ${T_{z}}$ инжектируемого сигнала (${T^{(k)}_{z}}$, ${k \in \left[1;K\right]}$):
\begin{itemize}
    \item Длительность ${T_h}$ измерительного окна равна периоду ${T_z}$ инжектируемого сигнала (${T^{(k)}_h = T^{(k)}_z}$, ${T^{(k)}_h - var}$).
    \item Период дискретизации ${T_s}$ одинаков для всех значений ${T_z}$ и равен периоду ${T_d}$ выполнения операций в системе (${T^{(k)}_s = T_d}$, ${T^{(k)}_s = const}$).
    \item Количество точек $N$ дискретной последовательности, по которым вычисляются коэффициенты разложения в выражении \eqref{eqn:DFT}, изменяется пропорционально ${T_z}$ (${N^{(k)} - var}$).
\end{itemize}

Период инжектируемого сигнала
\begin{equation*}
    T^{(k)}_z = N^{(k)} \mathord{\cdot} T^{(k)}_s = N^{(k)} \mathord{\cdot} T_d
\end{equation*}
кратен периоду ${T_d}$ выполнения операций в системе (рисунок~\ref{fig:N_var}).

\begin{figure}[H]
	\centering
	\includegraphics [width=1.00\linewidth] {frequencyResponse/N_var.drawio.png}

	\caption{Размер ${N^{(k)}}$ дискретной последовательности при изменении периода ${T^{(k)}_z}$ инжектируемого сигнала}
	\label{fig:N_var}
\end{figure}

В связи с тем, что ${T^{(k)}_h = T^{(k)}_z}$ составляющей с частотой ${f^{(k)}_z}$ соответствуют коэффициенты ${X^{(k)}\left[l\right]}$, ${Y^{(k)}\left[l\right]}$ всегда при ${l = 1}$, т.е. коэффициент передачи динамического звена на частоте ${f^{(k)}_z}$ определяется выражением \eqref{eqn:complexFreqResp} при ${L^{(k)} = 1}$.

Минимальное количество отсчетов $N$ на измерительном окне ${T_h}$ принято равным 3, поэтому частота ${f_z}$ ограничена сверху значением ${f_d / 3}$.

Таким образом, изложенному алгоритму измерения частотных характеристик с последовательной сменой частоты ${f_z}$ инжектируемого сигнала в серии из $K$ измерений \textbf{характерны следующие особенности}:
\begin{enumerate}
    \item Минимально возможная частота ${f_z}$ напрямую определяется максимальным значением измерительного окна ${T_h}$, которое ограничено сверху только объемом памяти для хранения массивов-откликов ${x\left[n\right]}$, ${y\left[n\right]}$.
    
    \item Максимально возможная частота ${f_z}$ в пределе составляет ${f_d / 3}$.
    
    \item В памяти достаточно хранить ${N^{(k)} = f_d / f^{(k)}_z}$-значений каждого массива-отклика ${x^{(k)}\left[n\right]}$, ${y^{(k)}\left[n\right]}$, причем величина $N$ не является постоянной и уменьшается с ростом частоты ${f_z}$ вплоть до ${N = 3}$ при ${f_z \rightarrow f_d / 3}$.
    
    \item Шаг приращения периода ${T_z}$ инжектируемого сигнала (разрешающая способность): ${\Delta T_z = T_d}$.
    
    \item Потенциальные значения частот ${f_z}$ инжектируемого сигнала неравномерно распределены на интервале ${f \in \left(0; {f_d / 2} \right)}$ возможных значений:
    \begin{itemize}
        \item На полуинтервале ${f \in \left[{f_d / 10}; {f_d / 2} \right)}$ частота ${f_z}$ может принимать значения всего из 8-ми возможных: \\
        ${f_z = \left({f_d / 10}, {f_d / 9}, ..., {f_d / 3} \right)}$.
        \item На полуинтервале ${f \in \left[{f_d / 100}; {f_d / 10} \right)}$ -- одно из 90 возможных: \\
        ${f_z = \left({f_d / 100}, {f_d / 99}, ..., {f_d / 11} \right)}$.
        \item На интервале ${f \in \left(0; {f_d / 100} \right)}$ -- сколь угодно много значений, ограниченное исключительно объемом памяти для хранения массивов-откликов ${x\left[n\right]}$, ${y\left[n\right]}$.
    \end{itemize}
\end{enumerate}
\section{Алгоритм с изменяемой длительностью измерительного окна, равной периоду инжектируемого сигнала, и децимацией}\label{sec:N_const}

\needspace{1\baselineskip}
Рассматриваемый алгоритм представляет собой модификацию алгоритма \ref{sec:N_var} в части периода дискретизации ${T_s}$ -- период дискретизации пропорционален периоду инжектируемого сигнала ${T_z}$. Ввиду плавающего ${T_s}$ длина $N$ последовательностей ${x\left[n\right]}$, ${y\left[n\right]}$, по которым вычисляются коэффициенты разложения в выражении \eqref{eqn:DFT}, неизменна для всех значений периода инжектируемого сигнала ${T_z}$. Условие кратности периоду выполнения операций в системе ${T_d}$ остается.

Постоянство $N$ при изменении периода инжектируемого сигнала ${T_z}$ обеспечивается децимацией (прореживанием) последовательностей ${x'[n]}$, ${y'[n]}$ с коэффициентом децимации ${M^{(k)} \in \mathbb{N}}$:
\begin{equation*}
    \begin{split}
        &x^{(k)}[n] = x'^{(k)}\left[n \mathord{\cdot} M^{(k)}\right] = x^{(k)}_c \left(n \mathord{\cdot} M^{(k)} \mathord{\cdot} T_d \right), \\
        &y^{(k)}[n] = y'^{(k)}\left[n \mathord{\cdot} M^{(k)}\right] = y^{(k)}_c \left(n \mathord{\cdot} M^{(k)} \mathord{\cdot} T_d \right),
    \end{split}
\end{equation*}
Переменными ${x'[n]}$, ${y'[n]}$ обозначены последовательности, формируемые в результате дискретизации откликов ${x_c(t)}$, ${y_c(t)}$ непосредственно с периодом выполнения операций в системе ${T_d}$.

Таким образом, в серии из $K$ измерений для каждого значения периода инжектируемого сигнала ${T^{(k)}_z}$, ${k \in \{1, \dots, K\}}$ (рисунок~\ref{fig:N_const}), выполняется:
\begin{equation*}
    T^{(k)}_h = T^{(k)}_z = N^{(k)} \mathord{\cdot} \underbrace{M^{(k)} \mathord{\cdot} T_d}_{T^{(k)}_s}.
\end{equation*}

\begin{figure}[H]
    \centering
    \includegraphics [width=1.00\linewidth] {frequencyResponse/N_const.pdf}

    \caption{Параметры в алгоритме с изменяемой длительностью измерительного окна, равной периоду инжектируемого сигнала, и децимацией}
    \label{fig:N_const}
\end{figure}

В связи с тем, что \textbf{так же, как и в алгоритме \ref{sec:N_var}}, длительность измерительного окна ${T^{(k)}_h}$ равна периоду инжектируемого сигнала ${T^{(k)}_z}$, составляющей с частотой ${f^{(k)}_z}$ соответствуют коэффициенты ${X^{(k)}\left[l\right]}$, ${Y^{(k)}\left[l\right]}$ всегда при ${l = 1}$, т.е. коэффициент передачи динамического звена на частоте ${f^{(k)}_z}$ определяется выражением \eqref{eqn:complexFreqResp} при ${L^{(k)} = 1}$. \textbf{Принципиальное отличие от алгоритма \ref{sec:N_var}} связано с потенциально меньшей длиной $N$ массивов-откликов вследствие децимации.

\needspace{1\baselineskip}
Для изложенного алгоритма измерения характерно следующее:
\begin{enumerate}
    \item Максимально возможный период инжектируемого сигнала ${T_z}$ соответствует максимальному значению измерительного окна ${T_h}$, но уже не ограничен параметром ${N_\text{max}}$, т.к. длина $N$ каждого массива-отклика ${x\left[n\right]}$, ${y\left[n\right]}$ неизменна (не зависит от периода инжектируемого сигнала ${T_z}$).
    \item Шаг приращения периода инжектируемого сигнала ${T_z}$ (разрешающая способность) в $N$ раз больше по сравнению с алгоритмом \ref{sec:N_var} (${\Delta T_z = N \mathord{\cdot} T_d}$).
    \item Грубый шаг приращения периода обусловливает еще более неравномерное распределение частот. Например, в диапазоне частот ${f \in \left[{f_d / 100}; {f_d / 10} \right]}$ при ${N = 10}$ частота инжектируемого сигнала ${f_z}$ может принимать значения всего из 10 возможных: \\
    ${f_z = \left({f_d / 100}, {f_d / 90}, ..., {f_d / 10} \right)}$.
\end{enumerate}

\section{Алгоритм с неизменной длительностью измерительного окна, кратной периоду инжектируемого сигнала}\label{sec:Th_const}

\needspace{1\baselineskip}
Принципиальные особенности рассматриваемого алгоритма:
\begin{itemize}
    \item Длительность измерительного окна ${T_h}$ неизменна и равна наибольшему среди $K$-значений периоду инжектируемого сигнала.
    \item Длина $N$ последовательностей ${x\left[n\right]}$, ${y\left[n\right]}$, по которым вычисляются коэффициенты разложения в выражении \eqref{eqn:DFT}, неизменна для всех значений периода инжектируемого сигала ${T^{(k)}_z}$.
    \item Период дискретизации ${T_s}$ одинаков для всех значений периода инжектируемого сигнала ${T_z}$ и кратен периоду выполнения операций в системе ${T_d}$.
\end{itemize}

Математическая обработка в соответствии с \eqref{eqn:DFT} накладывает ограничение на период ${T_z}$ инжектируемого сигнала -- требуется, чтобы измерительное окно ${T_h}$ вмещало целое число ${L^{(k)} \in \mathbb{N}}$ периодов ${T_z}$.

Таким образом, в серии из $K$ изменений для каждого значения периода инжектируемого сигнала ${T^{(k)}_z}$, ${k \in \{1, \dots, K\}}$ (рисунок~\ref{fig:Th_const}), выполняется:
\begin{equation*}
    T^{(k)}_h = L^{(k)} \mathord{\cdot} T^{(k)}_z = N^{(k)} \mathord{\cdot} \underbrace{M^{(k)} \mathord{\cdot} T_d}_{T^{(k)}_s},
\end{equation*}
причем ${T^{(k)}_h = const}$.

\begin{figure}[H]
    \centering
    \includegraphics [width=1.00\linewidth] {frequencyResponse/Th_const.pdf}

    \caption{Параметры в алгоритме с неизменной длительностью измерительного окна, кратной периоду инжектируемого сигнала}
    \label{fig:Th_const}
\end{figure}

Частота инжектируемого сигнала потенциально принимает значения
\begin{equation*}
    f^{(k)}_z = L^{(k)} \cfrac{1}{T_h} = L^{(k)} \Delta f_z,
\end{equation*}
где ${\Delta f_z}$ -- шаг приращения частоты инжектируемого сигнала. Составляющей с частотой ${f^{(k)}_z}$ соответствуют коэффициенты ${X^{(k)}\left[l\right]}$, ${Y^{(k)}\left[l\right]}$ выражения \eqref{eqn:complexFreqResp} при ${l = L^{(k)}}$.

В связи с ${\Delta f_z = const}$ потенциальные значения частот инжектируемого сигнала ${f_z}$ равномерно распределены на интервале возможных значений, что \textbf{является главной отличительной особенностью алгоритма} в сравнении с алгоритмами \ref{sec:N_var}, \ref{sec:N_const}.

\needspace{1\baselineskip}
Изложенному алгоритму измерения характерно следующее:
\begin{enumerate}
    \item Максимально возможный период инжектируемого сигнала ${T_z}$ ограничен только параметром ${N_\text{max}}$.

    \item Длина $N$ каждого массива-отклика ${x\left[n\right]}$, ${y\left[n\right]}$ неизменна (не зависит от периода инжектируемого сигнала ${T_z}$).
    
    \item В случае ${N_\text{min}=3}$ максимально возможная частота ${f_z}$ ограничена сверху значением ${f_d / 3}$.

    \item Шаг приращения частоты инжектируемого сигнала ${f_z}$ (разрешающая способность): ${\Delta f_z = 1 / T_h}$.
    
    \item Период $T_z$ инжектируемого сигнала потенциально принимает значения от ${T^{min}_z = 3 T_s}$ до ${T^{max}_z = T_h}$, но при условии ${T_h / T_s \in \mathbb{N}}$.
\end{enumerate}
\section{Алгоритм с изменяемой длительностью измерительного окна, кратной периоду инжектируемого сигнала}\label{sec:fr_optim}

\needspace{1\baselineskip}
{Принципиальные отличия рассматриваемого алгоритма от алгоритма \ref{sec:Th_const}} заключаются в том, что длительность измерительного окна ${T_h}$, длина последовательностей $N$ и период дискретизации ${T_s}$ -- плавающие параметры, соответствующие каждому отдельному значению периода инжектируемого сигнала ${T_z}$. Остаются следующие ограничения:
\begin{itemize}
    \item Возможные значения длительности измерительного окна ${T_h}$ кратны периоду инжектируемого сигнала ${T_z}$.

    \item Период дискретизации ${T_s}$ кратен периоду выполнения операций в системе ${T_d}$.
\end{itemize}
Эти ограничения сводятся к плавающим параметрам ${L, M \in \mathbb{N}_{>0}}$, в дополнение к плавающему параметру ${N \in \mathbb{N}_{>0}}$.

Таким образом, в серии из $K$ измерений для каждого значения периода инжектируемого сигнала ${T^{(k)}_z}$, ${k \in \{1, \dots, K\}}$ (рисунок~\ref{fig:Th_var_N_var}), выполняется:
\begin{equation}\label{eqn:windowlength}
    T^{(k)}_h = L^{(k)} \mathord{\cdot} T^{(k)}_z = N^{(k)} \mathord{\cdot} \underbrace{M^{(k)} \mathord{\cdot} T_d}_{T^{(k)}_s} = L^{(k)} \mathord{\cdot} R^{(k)} \mathord{\cdot} T_d,
\end{equation}
где ${R^{(k)} \in \mathbb{R}_{>0}}$ -- отношение периода инжектируемого сигнала ${T^{(k)}_z}$ к периоду выполнения операций в системе ${T_d}$.

\begin{figure}[H]
    \centering
    \includegraphics [width=1.00\linewidth] {frequencyResponse/Th_var_N_var.pdf}

    \caption{Параметры в алгоритме с изменяемой длительностью измерительного окна, кратной периоду инжектируемого сигнала}
    \label{fig:Th_var_N_var}
\end{figure}

\textbf{Аналогично алгоритму \ref{sec:Th_const}}, соотношение между периодом инжектируемого сигнала и периодом дискретизации может принимать дробные значения. В совокупности с изменяемой длительностью измерительного окна это обеспечивает на интервале ${f \in \left(0; {f_d / 2} \right)}$ сколь угодно много потенциальных значений частот ${f_z}$ инжектируемого сигнала, что \textbf{является главной отличительной особенностью алгоритма}.

Из равенства \eqref{eqn:windowlength}
\begin{equation}\label{eqn:L_ratio}
    L^{(k)} = \frac{N^{(k)} \mathord{\cdot} M^{(k)}}{R^{(k)}}.
\end{equation}

В связи с тем, что частота инжектируемого сигнала кратна частоте первой гармоники в спектре:
\begin{equation*}
    f^{(k)}_z = L^{(k)} \mathord{\cdot} f^{(k)}_h = \cfrac{L^{(k)}}{T^{(k)}_h},
\end{equation*}
и ограничена половиной частоты дискретизации:
\begin{equation*}
    f^{(k)}_z < \cfrac{f^{(k)}_s}{2} \quad \rightarrow \quad \cfrac{L^{(k)}}{T^{(k)}_h} < \cfrac{f_d}{2 M^{(k)}},
\end{equation*}
множитель ${L^{(k)}}$ ограничен диапазоном
\begin{equation*}
    1 \leq L^{(k)} < N^{(k)} / 2,
\end{equation*}
что, с учетом \eqref{eqn:L_ratio}, накладывает ограничение на коэффициент децимации
\begin{equation*}
    1 \leq M^{(k)} < R^{(k)} / 2.
\end{equation*}

Таким образом, для каждой $k$-й частоты ${f_z}$ задача сводится к поиску наборов ${\mathbf{s}_{(p)} \in \mathbb{N}_{> 0}^3}$ натуральных чисел ${N_{(p)}, M_{(p)}, L_{(p)}}$, при которых частота инжектируемого сигнала ${f_{z(p)}}$
\begin{equation*}
    f_{z(p)} = \cfrac{L_{(p)}}{N_{(p)} \mathord{\cdot} M_{(p)}} f_d
\end{equation*}
отклоняется от желаемой частоты ${f_z}$ в пределах заданного значения ${\Delta\epsilon^{*}_{f}}$:
\begin{equation*}
    \left| f_{z(p)} - f_z \right| = \Delta\epsilon_{f(p)} \leq \Delta\epsilon^{*}_f,
\end{equation*}
при ограничениях
\begin{equation*}
        N_\text{min} \leq N_{(p)} \leq N_\text{max}, \quad 1 \leq M_{(p)} < R/2.
\end{equation*}

Оптимальным считается такой набор ${\mathbf{s}_{(p^{*})}}$, в котором параметр ${N^{*}}$ минимален среди всех наборов ${\mathbf{s}_{(p)}}$, отвечающих условию
\begin{equation}\label{eqn:freqtol}
    \Delta\epsilon_{f(p)} \leq \Delta\epsilon^{*}_{f}.
\end{equation}
В случае, когда ни один набор параметров ${\mathbf{s}_{(p)}}$ не удовлетворяет условию \eqref{eqn:freqtol}, оптимальным считается набор, соответствующий минимальному отклонению ${\Delta\epsilon_{f(p)}}$. Таким образом, согласно выражению \eqref{eqn:meas_duration}, минимум длины массивов-откликов $N$ обеспечивает сокращение продолжительности измерения частотной характеристики.
\section{Алгоритм поиска оптимального набора параметров}

Поиск оптимального набора параметров требуется для каждой частоты инжектируемого сигнала из исходного набора частот ${f^{(k)}_{z}}$, ${k \in \left[1;K \right]}$. В результате определяются такие $K$-наборов параметров, при котором частотная характеристика измеряется алгоритмом \ref{sec:fr_optim} за минимальное время с точностью воспроизведения частоты ${f^{(k)}_{z}}$ близкой к заданной ${\Delta\epsilon^{*}_{f}}$.

\textbf{Входные данные:}
\begin{itemize}
    \item Частота инжектируемого сигнала: ${f_z}$.
    \item Максимальное абсолютное отклонение от заданной частоты ${f_z}$ инжектируемого сигнала: ${\Delta\epsilon^{*}_{f}}$.
    \item Частота выполнения операций в системе: ${f_d}$.
    \item Максимальное количество точек дискретных последовательностей ${x[n]}$, ${y[n]}$ на интервале измерения: ${N_\text{max}}$.
    \item Минимальное количество точек дискретных последовательностей ${x[n]}$, ${y[n]}$ на интервале измерения: ${N_\text{min}}$.
\end{itemize}

\textbf{Выходные данные:}
\begin{itemize}
    \item Количество точек дискретных последовательностей ${x[n]}$, ${y[n]}$ на интервале измерения: ${N}$.
    \item Коэффициент децимации: ${M}$.
    \item Коэффициент кратности длительности измерительного окна ${T_h}$ к периоду инжектируемого сигнала ${T_z}$: ${L}$.
\end{itemize}

\textbf{Последовательность действий:}
\begin{enumerate}
    \item Вычисляется коэффициент соотношения между периодом инжектируемого сигнала и периодом выполнения операций в системе:
        \begin{equation*}
            R = \frac{f_{d}}{f_{z}}.
        \end{equation*}
    \item Вычисляется максимально возможное значение коэффициента децимации:
        \begin{equation*}
            M_\text{max} = \left\lfloor \frac{R}{2} \right\rfloor.
        \end{equation*}
    \item В цикле ${M \in \left[ 1; M_\text{max} \right]}$:
    \begin{enumerate}
        \item Вычисляются всевозможные значения
            \begin{equation*}
                L = \frac{M \cdot N}{R}
            \end{equation*}
        для ${N \in \left[ N_\text{min}; N_\text{max} \right]}$, с округлением $L$ до 1 при ${L<1}$.
        \item Значения $L$ используются для вычисления всевозможных приближений к начальному значению коэффициента ${R}$:
            \begin{equation*}
                R^{'} = \frac{M \cdot N}{\left\lfloor L \right\rfloor}.
            \end{equation*}
        \item Значения ${R^{'}}$ используются для вычисления всевозможных приближений к исходному значению частоты ${f_{z}}$:
            \begin{equation*}
                f^{'}_{z} = \frac{f_d}{R^{'}},
            \end{equation*}
            и соответствующих отклонений
            \begin{equation*}
                \Delta\epsilon_{f} = \left| f^{'}_{z} - f_{z} \right|.
            \end{equation*}
        \item Из набора значений ${f^{'}_{z}}$ исключаются те, для которых ${\Delta\epsilon_{f} > \Delta\epsilon^{*}_{f}}$.
        \item Среди оставшихся ${f^{'}_{z}}$ определяется значение, соответствующее минимальному ${N}$. Формируется набор ${\mathbf{s}^{*}}$ соответствующих параметров ${N, M, L}$.
        \item Если нет значений ${f^{'}_{z}}$, удовлетворяющих условию ${\epsilon_{f} < \epsilon^{*}_{f}}$, выбирается набор параметров ${N, M, L}$, соответствующий наименьшему отклонению ${\Delta\epsilon_{f}}$.
    \end{enumerate}
    \item Среди найденных наборов параметров ${\mathbf{s}^{*}}$, соответствующих разным $M$, рассматриваются наборы, удовлетворяющие условию ${\Delta\epsilon_{f} < \Delta\epsilon^{*}_{f}}$ и среди них выбирается набор с минимальным ${N}$.
    \item Если нет наборов параметров ${\mathbf{s}^{*}}$, удовлетворяющих условию ${\Delta\epsilon_{f} < \Delta\epsilon^{*}_{f}}$, выбирается набор параметров, соответствующий наименьшему отклонению ${\Delta\epsilon_{f}}$.
    \item Формируется итоговый набор ${\mathbf{s}^{*} = \left( N, M, L \right)}$.
\end{enumerate}

\textbf{Примечания:}
\begin{enumerate}
    \item Отклонение от заданной частоты ${f_z}$ инжектируемого сигнала допустимо задавать в относительных единицах ${\epsilon^{*}_{f}}$, с последующим пересчетом в абсолютное значение
    \begin{equation*}
        \Delta\epsilon^{*}_{f} = \epsilon^{*}_{f} \cdot f_z.
    \end{equation*}
    В таком случае, на вход алгоритма получается задать общее для всех частот относительное отклонение ${\epsilon^{*}_{f}}$. Абсолютные значения отклонения ${\Delta\epsilon^{*(k)}_{f}}$ сформируются в соответствии частотами ${f^{(k)}_z}$.
    \item Максимальное количество точек ${N_\text{max}}$ определяется размером памяти, доступной в буфере встраиваемой системы управления, зарезервированном под внутрисхемное измерение частотных характеристик.
\end{enumerate}


\textbf{Пример.}

Заданы: ${f_d=100}$~кГц, ${N_\text{min}=10}$, ${N_\text{max}=1214}$.

Определить набор параметров ${N, M, L}$ для измерения на инжектируемой частоте ${f_z=22,1}$~кГц с относительным отклонением до 0,1\% ${\left(\epsilon^{*}_{f}=0,001\right)}$.

Коэффициент соотношения между периодом инжектируемого сигнала и периодом выполнения операций в системе:
\begin{equation*}
    R = \frac{f_{d}}{f_{z}} = \frac{100000}{22100} = 4,524886877828054.
\end{equation*}

Максимально возможное значение коэффициента децимации:
\begin{equation*}
    M_\text{max} = \left\lfloor \frac{R}{2} \right\rfloor = 2.
\end{equation*}

Оптимальный набор параметров: ${N=43}$, ${M=2}$, ${L=19}$.

Соответствующее приближение к начальному значению коэффициента ${R}$:
\begin{equation*}
    R^{'} = \frac{M \cdot N}{\left\lfloor L \right\rfloor} = \frac{2 \cdot 43}{19} = 4,526315789473684.
\end{equation*}

Соответствующее приближение к исходному значению частоты ${f_{z}}$:
\begin{equation*}
    f^{'}_{z} = \frac{L}{M \cdot N}f_d = \frac{19}{2 \cdot 43} = 22093,023255813954~\text{Гц}.
\end{equation*}

Относительное отклонение частоты
\begin{equation*}
    \epsilon_{f} = \frac{\left| f^{'}_{z} - f_{z} \right|}{f_{z}} = 0,0003156897821740197
\end{equation*}
не превышает заданного ограничения ${\epsilon^{*}_{f}=0,001}$.

Таким образом, чтобы воспроизвести частоту инжектируемого сигнала с точностью до 0,1\% достаточно измерить ${N=43}$ точки на ${L=19}$ периодах инжектируемого сигнала. Период выборки при этом должен быть в ${M=2}$ раза больше периода выполнения операций в системе.
\section{Результаты измерений}

Оценка качественных характеристик разработанного алгоритма внутрисхемного измерения частотных характеристик выполнена в серии экспериментов по измерению частотных характеристик динамических звеньев системы автоматического регулирования источника питания (Приложение~\ref{appendix:bench}). Частотные характеристики измерены в диапазоне частот от 10~Гц до 49~кГц. Измерения выполнены двумя способами: сертифицированным измерителем частотных характеристик (AP300) и разработанным алгоритмом внутрисхемного измерения \ref{sec:fr_optim} (Digital Points).

\begin{figure}[H]
	\centering
	\includegraphics [width=1.00\linewidth] {frequencyResponse/Wui.pdf}

	\caption{Результаты эксперимента}
	\label{fig:Wui_FR}
\end{figure}

\begin{figure}[H]
	\centering
	\includegraphics [width=1.00\linewidth] {frequencyResponse/Wui_delay.pdf}

	\caption{Результаты эксперимента}
	\label{fig:Wui_delay_FR}
\end{figure}

\begin{figure}[H]
	\centering
	\includegraphics [width=1.00\linewidth] {frequencyResponse/G_I_PI.pdf}

	\caption{Результаты эксперимента}
	\label{fig:G_I_PI_FR}
\end{figure}

\begin{figure}[H]
	\centering
	\includegraphics [width=1.00\linewidth] {frequencyResponse/T_I_PI.pdf}

	\caption{Результаты эксперимента}
	\label{fig:T_I_PI_FR}
\end{figure}

\begin{figure}[H]
	\centering
	\includegraphics [width=1.00\linewidth] {frequencyResponse/G_I_PID.pdf}

	\caption{Результаты эксперимента}
	\label{fig:G_I_PID_FR}
\end{figure}

\begin{figure}[H]
	\centering
	\includegraphics [width=1.00\linewidth] {frequencyResponse/G_U_PI.pdf}

	\caption{Результаты эксперимента}
	\label{fig:G_U_PI_FR}
\end{figure}